\subsection{Scope and Remaining Validation}
\label{sec:design-discussion}

\paragraph{Slack error and profile drift.}
Overestimated slack can place a critical request in Tier R, forcing
history recomputation after failure. Underestimated slack admits more
requests to Peer competition without changing their footprints.
P90 estimation, hysteresis, unknown-successor handling, and event-driven
replanning mitigate these errors but do not bound them. An opportunistic
Tier-R DRAM copy helps only after it has committed.
Hot-set identities are stale predictions from the previous ranking
window (default 512 decoding steps), whereas union sizing uses a window
matched to replay length (default 64 queries). Slow profile refreshes
can track drift, but do not guarantee that admission never changes
discontinuously. Slack error, ranking-window length, and profile age
require evaluation under agent workloads.

\paragraph{Selector and execution scope.}
Integration requires actual selection indices, recoverable selector
state, and all required attention branches. Token selectors must be
mapped to physical copy blocks, potentially amplifying transfer bytes;
compression and local-window branches cannot be omitted.
Unstable rankings or a poor granularity match can fall back to full-layer
loading. For dense attention, $h=\widetilde u=1$, and only layer order
remains. Standalone requests have $\sigma_r=0$ and use stable IDs for ties;
their recovery is still concurrent. A failure during prefill may lose
an entire in-flight chunk, requiring a replay-length-specific calibration
instead of the 64-token decode-tail assumption.

\paragraph{Evidence and limitations.}
Timing prediction by role remains unvalidated: the available CC
labels collapse all agents into one role, and the workflow-disjoint
diagnostic does not demonstrate accurate timing. The sparse-model survey
documents representative adoption rather than a deployment share.
Retry-demand coverage, access concentration, and union statistics come from native selector
traces on three book documents per model. They need not transfer to
code and tool-heavy agent histories. The retry-demand experiment predicts
a held-out causal suffix; it does not inject a failure, restore selector
state, or measure replay time. Larger-context footprints extrapolate
32K union fractions; the NOSA profile comparison borrows MiniCPM4's
compute and layer-byte configuration. The pipeline and footprint figures
are calibrated models, not physical failure-recovery measurements.
End-to-end validation of the new placement policy, simultaneous retries,
foreground interference, and runtime overhead remains outstanding.
Older serial-recovery or previous-policy results do not establish these
properties. The invariants above specify intended runtime behavior.
